% Environnements personnalisés LaTeX
% Fichier: config/environments.tex

% ============================================================================
% ENVIRONNEMENT POUR LES VULNÉRABILITÉS
% ============================================================================

\newenvironment{vulnerability}[3]{%
    % #1 = Nom de la vulnérabilité
    % #2 = Catégorie OWASP
    % #3 = Niveau de difficulté
    \begin{tcolorbox}[
        colback=primarycolor!5!white,
        colframe=primarycolor!75!black,
        title=\textbf{#1},
        fonttitle=\bfseries,
        subtitle style={colback=primarycolor!10!white},
        subtitle={\textbf{Catégorie OWASP:} #2 \hfill \textbf{Difficulté:} #3},
        breakable
    ]
}{%
    \end{tcolorbox}
}

% ============================================================================
% ENVIRONNEMENT POUR LES PAYLOADS D'EXPLOITATION
% ============================================================================

\newenvironment{payload}{%
    \begin{tcolorbox}[
        colback=dangercolor!5!white,
        colframe=dangercolor!75!black,
        title=\textbf{Payload d'exploitation},
        fonttitle=\bfseries\small,
        breakable
    ]
    \begin{verbatim}
}{%
    \end{verbatim}
    \end{tcolorbox}
}

% ============================================================================
% ENVIRONNEMENT POUR LES RECOMMANDATIONS
% ============================================================================

\newenvironment{recommendation}[1]{%
    % #1 = Priorité (Critique, Élevée, Moyenne, Basse)
    \begin{tcolorbox}[
        colback=successcolor!5!white,
        colframe=successcolor!75!black,
        title=\textbf{Recommandation (#1)},
        fonttitle=\bfseries\small,
        breakable
    ]
}{%
    \end{tcolorbox}
}

% ============================================================================
% ENVIRONNEMENT POUR LES ALERTES DE SÉCURITÉ
% ============================================================================

\newenvironment{securityalert}[1]{%
    % #1 = Niveau d'alerte
    \begin{tcolorbox}[
        colback=warningcolor!10!white,
        colframe=warningcolor!75!black,
        title=\textbf{Alerte de Sécurité: #1},
        fonttitle=\bfseries,
        arc=0mm,
        boxrule=0.5pt,
        breakable
    ]
}{%
    \end{tcolorbox}
}

% ============================================================================
% ENVIRONNEMENT POUR LE CODE SOURCE
% ============================================================================

\newenvironment{code}[1]{%
    % #1 = Langage (sql, js, html, python, etc.)
    \lstset{language=#1}
    \begin{lstlisting}
}{%
    \end{lstlisting}
}

% ============================================================================
% ENVIRONNEMENT POUR LES COMMANDES TERMINAL
% ============================================================================

\newenvironment{terminal}{%
    \begin{tcolorbox}[
        colback=darkcolor!95!white,
        colframe=darkcolor,
        title=\textbf{Terminal},
        fonttitle=\bfseries\small\color{lightcolor},
        coltext=lightcolor,
        breakable
    ]
    \begin{verbatim}
}{%
    \end{verbatim}
    \end{tcolorbox}
}

% ============================================================================
% ENVIRONNEMENT POUR LES NOTES PRATIQUES
% ============================================================================

\newenvironment{practicenote}{%
    \begin{tcolorbox}[
        colback=infocolor!5!white,
        colframe=infocolor!75!black,
        title=\textbf{Note Pratique},
        fonttitle=\bfseries\small,
        breakable
    ]
}{%
    \end{tcolorbox}
}

% Fin du fichier environments.tex
